% ch02d —— 原書 p.77–86（§2.3 末、§2.4、§2.5 前段）
% 命题 2.42 的陈述止于 p.76（由 ch02c 写完），本块从 p.77 的证明开始。

% [原文 p.77]
\begin{proof}
设 $Z$ 为定义 2.37 中 dlt 性所涉及的闭子集。若 $x\notin Z$，则化为第 1 或第 2 种情形。以下设 $x\in Z$，证明化为第 2 种情形。设 $E\subset Y$ 为双有理态射 $f:Y\to X$ 的例外除子。若 $f(E)\subset Z$，则由 dlt 的定义，$a(E,X,\Delta)>-1$ 成立。若 $f(E)\not\subset Z$，则 $\Delta$ 的两个不可约分支在 $x$ 的穿孔邻域内不相交，故由推论 2.31(3) 得 $a(E,X,\Delta)\ge 0$。因此化为第 2 种情形。
\end{proof}

\begin{propositionp}{2.43}
设 $(X,\Delta)$ 为 dlt 对，$X$ 为具有丰富除子 $H$ 的射影簇。设 $\Delta_1$ 是使 $\Delta-\Delta_1$ 为有效的有效 $\Q$-除子（不必为 $\Q$-Cartier 除子）。此时存在有理数 $c>0$ 与有效 $\Q$-除子 $D$，它们满足：$D$ 与 $\Supp\Delta$ 无公共分支，且 $D-\Delta_1-cH\equiv 0$；并且对一切有理数 $0<\varepsilon\ll 1$，$(X,\Delta-\varepsilon\Delta_1+\varepsilon D)$ 为 dlt。

特别地，当 $\Supp\Delta_1=\Supp\Delta$ 时，$(X,\Delta-\varepsilon\Delta_1+\varepsilon D)$ 为 klt。
\end{propositionp}

\begin{proof}
取 $m,m'\in\Z_{>0}$，使得 $m\Delta_1$ 为除子，且层 $\cO_X(m\Delta_1+m'H)$ 由全局截面生成。取一般元 $D'\in|m\Delta_1+m'H|$，并记 $D=(1/m)D'$。令 $c=m'/m$，则 $D-\Delta_1-cH\equiv 0$。由于 $-m\Delta_1+mD\sim m'H$ 为 Cartier 除子，故 $K_X+\Delta-\varepsilon\Delta_1+\varepsilon D$ 为 $\Q$-Cartier 除子。

取定义 dlt 性时出现的 $Z\subset X$。$m\Delta_1$ 在 $X\setminus Z$ 上为 Cartier 除子，因此 $|m\Delta_1+m'H|$ 在 $X\setminus Z$ 上无基点。由于 $D'$ 为一般元，$\Delta+D$ 在 $X\setminus Z$ 上是简单正规交叉除子。

设 $f:Y\to X$ 为 $(X,\Delta+D)$ 的任一对数奇点消解。对任一不可约除子 $E\subset f^{-1}(Z)$，有
\[
a(E,X,\Delta-\varepsilon\Delta_1+\varepsilon D)\longrightarrow a(E,X,\Delta)>-1\quad(\varepsilon\to 0).
\]
故由命题 2.40，$(X,\Delta-\varepsilon\Delta_1+\varepsilon D)$ 对 $0<\varepsilon\ll 1$ 为 dlt。最后一项断言由命题 2.41 推出。
\end{proof}

% [Tutor Note: 结论中的 $D$ 必须避开 $\Supp\Delta$ 的分支，否则 $\Delta+\varepsilon(D-\Delta_1)$ 会在这些分支上叠加出大于 1 的系数，dlt 立刻失效；$\Delta_1$ 的系数则被按 $1-\varepsilon$ 缩向 0，这正是把 dlt 对扰动成 klt 对的手法。]

\begin{theorem}[{[Sza95]}]{2.44}
设 $(X,\Delta=\sum d_iD_i)$，其中 $D_i$ 为互不相同的素除子，系数 $d_i$ 满足 $0\le d_i\le 1$。此时以下条件等价：
% [原文 p.78]
\begin{enumerate}[label=（\arabic*）,leftmargin=2.8em]
\item $(X,\Delta)$ 为 dlt。
\item 存在 $(X,\Delta)$ 的对数奇点消解 $f:X'\to X$，使得所有例外除子 $E_i\subset X'$ 都满足 $a(E_i,X,\Delta)>-1$。
\end{enumerate}
\end{theorem}

最后，回顾 Zariski 的一个结果：任一除子都可由有限次爆破得到。许多情形下，这一结果可代替奇点消解。证明取自 [Art86]。

\begin{lemmap}{2.45}
设 $X$ 为域 $k$ 上的代数簇。任取 $X$ 的分式域 $k(X)$ 的一个离散赋值环 $(R,\mathfrak m_R)$，使其满足 $\operatorname{trdeg}(R/\mathfrak m_R:k)=\dim X-1$。由此自然确定概形 $Y=\Spec R$、闭点 $y\in Y$、双有理态射 $f:Y\to X$ 等。此时按如下方式定义簇与态射的序列：
\[
X_0=X,\qquad f_0=f.
\]
假设 $f_i:Y\to X_i$ 已定义。令 $Z_i\subset X_i$ 为 $x_i=f_i(y)$ 的闭包。以 $Z_i$ 为中心作爆破 $X_{i+1}=B_{Z_i}X_i$，并令 $f_{i+1}:Y\to X_{i+1}$ 为自然确定的态射。
则对适当的 $n\ge 0$，由 $f_n:Y\to X_n$ 导出同构 $\cO_{X_n,x_n}\cong R$。
\end{lemmap}

\begin{proof}
回忆离散赋值环 $R$ 通过下式定义赋值 $v_R$：
\[
v_R(g):=\max\{s\in\Z: g\in\mathfrak m_R^s\}\qquad(g\in k(X)\setminus\{0\}).
\]
把上述爆破序列用环论语言改写。首先给定局部环 $\cO_0:=\cO_{X_0,x_0}$ 与其极大理想 $\mathfrak m_0$。假设 $\cO_n$ 与 $\mathfrak m_n$ 已定义。此时选取 $\mathfrak m_n$ 的生成元 $z_1,\dots,z_r$，使 $v_R(z_1)\le\cdots\le v_R(z_r)$。令 $\cO'_n:=\cO_n[z_2/z_1,\dots,z_r/z_1]$，并令 $\cO_{n+1}$ 为 $\cO'_n$ 关于 $\cO'_n\cap\mathfrak m_R$ 的局部化，$\mathfrak m_{n+1}$ 为其极大理想。

下面证明 $R=\bigcup_n\cO_n$。为此任取 $u\in R\setminus\cO_n$，把 $u$ 写成 $u=y_1/y_2$，其中 $y_i\in\cO_n$ 的选取使 $v_R(y_2)\ge 0$ 为最小。于是由下式可知 $y_i\in\mathfrak m_n$：
\[
u\notin\cO_n\iff y_2\in\mathfrak m_n\iff v_R(y_2)>0\Rightarrow y_1\in\mathfrak m_n.
\]
% [原文 p.79]
从而可把 $y_1$ 写成 $y_1=\sum_j z_jy_{1j}$，其中 $y_{1j}\in\cO_n$，$z_j$ 为上面选定的 $\mathfrak m_n$ 的生成元。令 $y'_1:=\sum_j(z_j/z_1)y_{1j}\in\cO_{n+1}$。则 $y_1=z_1y'_1$，$u=y'_1/y'_2$。注意由构造有 $v_R(y'_2)<v_R(y_2)$。反复进行此步骤，最终有 $u=y_1^{(s)}/y_2^{(s)}$，$v_R(y_2^{(s)})=0$。这意味着对适当的 $s$ 有 $u\in\cO_{n+s}$。至此 $R=\bigcup_n\cO_n$ 证毕。

设 $u_1,\dots,u_r\in k(X)$ 为 $R/\mathfrak m_R$ 在 $k$ 上的超越基。此时对 $n\gg 1$ 有 $u_1,\dots,u_r\in\cO_n$。因此
\[
\operatorname{trdeg}(k(x_n):k)=\operatorname{trdeg}(R/\mathfrak m_R:k)=\dim X-1\qquad(n\gg 1).
\]
由构造，$\cO_n$ 是有限生成 $k$-代数的一个局部化，故
\[
\operatorname{trdeg}(k(X):k)-\operatorname{trdeg}(k(x_n):k)=\dim\cO_n
\]
成立（参看 [Mat69, 14.C]）。因此 $\dim\cO_n=1$（$n\gg 1$），且 $R$ 是 $\cO_n$ 的正规化 $R'$ 的局部化。理由如下：$R$ 支配 $R'$ 的某个局部化 $R''$，而 $R'$ 的任一局部化都是离散赋值环，故 $R=R''$。又因 $\cO_n$ 是域上有限生成代数的局部化，$R'$ 是有限生成 $\cO_n$-模（[Mat69, 31.H]）。取其在 $\cO_n$ 上的生成元 $v_1,\dots,v_s$。对 $p\gg n$ 有 $v_1,\dots,v_s\in\cO_p$，故 $p\gg 1$ 时 $\cO_p$ 支配上述 $R''$。这意味着 $\cO_p=R$。
\end{proof}

% [Tutor Note: 由 $u\notin\cO_n$ 推出 $y_2\in\mathfrak m_n$，再由 $v_R$ 的最小性得 $y_1\in\mathfrak m_n$；每一步把分母的赋值压低，正是“用爆破逼近任意赋值”的代数影子。]

\section{小平消没定理}

\takeaway{本节用 Hodge 理论把全纯上同调提升到拓扑上同调，再经 GAGA 原理与循环覆盖化归到 Serre 消没定理，证明丰富线丛 $L$ 满足 $H^i(X,L^{-1})=0$（$i<\dim X$）。}

本节证明小平消没定理。下一节讨论各种精细化的版本；这些更一般的消没定理在第 3 章的论证中本质地需要。证明基于下面的原理。

\begin{principle}{2.46}
若层 $F$ 的上同调来自拓扑上同调，则存在小平型消没定理。
\end{principle}

关于这一原理的详细说明见 [Kol86b, Section 5]。这里通过用它证明古典的消没定理来说明该原理。证明取自 [Kol95]。

% [原文 p.80]
\begin{theorem*}[小平消没定理]{2.47}
设 $X$ 为非奇异射影簇，$L$ 为其上的丰富线丛。此时成立
\[
H^i(X,L^{-1})=0,\qquad \forall i<\dim X.
\]
\end{theorem*}

\begin{proof}
\textbf{步骤 1（GAGA 原理）}\quad 设 $X$ 为 $\C$ 上的有限型概形。$X$ 也可看作复解析空间，此时记作 $X^{an}$。若 $F$ 是 $X$ 上的凝聚层，则 $F^{an}$ 表示与之伴随的解析凝聚层。就茎而言，这一解析化有如下表示：
\[
F^{an}_{x,X^{an}}=F_{x,X}\otimes_{\cO_{x,X}}\cO^{an}_{x,X^{an}}.
\]
所谓的 GAGA 原理（[Ser56] 论文名的首字母缩写）断言：许多情形下，代数的对象与解析的对象表现相同。下面的定理陈述的是高次顺像，但目前只需上同调方面的结果。

\begin{theorem}{2.48}
设 $X,Y$ 为 $\C$ 上的有限型分离概形，$f:X\to Y$ 为固有态射，$F$ 为 $X$ 上的凝聚层。此时成立
\[
(R^if_*F)^{an}\cong R^i(f^{an})_*(F^{an})\qquad\forall i.
\]
若 $X$ 是固有的，则 $H^i(X,F)\cong H^i(X^{an},F^{an})$，$\forall i$。
\end{theorem}

由这一原理，本节以下只需在复解析的框架中计算上同调。为简化记号，将省略上标 $an$。

\textbf{步骤 2}\quad 凝聚层 $F$ 也是 Abel 群的层。要找的是拓扑可构造层 $\mathbf{F}$ 与自然映射 $\mathbf{F}\to F$，使其诱导的上同调群之间的映射为满射。（解析凝聚层的层上同调与其作为 Abel 群的层的上同调相同。）对 $F=\cO_X$，常层 $\C_X$ 即满足条件。这是因为由 Hodge 理论（例如参看 [GH78, p.116]），由层的包含关系确定的自然映射 $H^i(X,\C_X)\to H^i(X,\cO_X)$ 是满射。

\textbf{步骤 3（循环覆盖）}\quad 循环覆盖的构造在许多不同场合都有用，故分 3 步给出一般定义。

\begin{definition}[不分歧循环覆盖]{2.49}
设 $X$ 为正规簇（或正规复解析流形），$L$ 为 $X$ 上的线丛，且对某个 $m>0$ 有 $L^m\cong\cO_X$。对应的循环覆盖是次数为 $m$ 的有限平展（étale）态射 $p:X_{m,L}\to X$，它可用 3 种等价方式定义：
% [原文 p.81]
\begin{enumerate}[label=（\arabic*）,leftmargin=2.8em]
\item 拓扑定义。由假设 $c_1(L)\in H^2(X,\Z)$ 是 $m$-挠元。因此它对应商映射 $\pi_1(X)\to H_1(X,\Z)\to\Z_m$。这说明它定义了 $X$ 的 $m$-重覆盖。
\item 几何定义。把 $L$ 看作 $\C^\times$-丛 $L\to X$。设 $h:L\to X\times\C$ 为由 $(x,l)\mapsto(x,l^{\otimes m})$ 确定的态射，记 $X_{m,L}:=h^{-1}(X\times\{1\})$。
\item 代数定义。$\bigoplus_{i=0}^{m-1}L^{-i}$ 是 $X$ 上的代数层，其乘法在 $i+j\ge m$ 时由
\[
L^{-i}\otimes L^{-j}\cong L^{-i-j}\cong L^{-i-j+m}
\]
定义。记 $X_{m,L}:=\Spec_X\bigoplus_{i=0}^{m-1}L^{-i}$。（指数取负只是一种偏好：正负皆可。若 $L^m\cong\cO_X$，则 $(L^{-1})^m\cong\cO_X$，两种取法给出同构的覆盖。这里的选择与下面考察的分歧覆盖的记号相协调。）
\end{enumerate}
由上述 3 种描述中的任何一种都可知，$p^*L$ 是 $X_{m,L}$ 上的平凡线丛。例如在代数上论证，可得
\[
p_*(p^*L)\cong L\otimes\bigoplus_{i=0}^{m-1}L^{-i}\cong\bigoplus_{i=0}^{m-1}L^{1-i}.
\]
取 $i=1$ 得直和项 $L^0\cong\cO_X$。$\cO_X$ 的常截面确定 $p^*L$ 的无处为零的截面。
\end{definition}

% [Tutor Note: 三种定义给出同一个覆盖，关键在 $L^m\cong\cO_X$ 使 $\bigoplus L^{-i}$ 成为有限代数层；用负指数是为了让下面分歧覆盖中的关系 $z^m=s$ 与代数定义直接吻合。]

\begin{definition}[分歧循环覆盖，线丛的情形]{2.50}
设 $X$ 为正规簇（或正规解析流形），$L$ 为 $X$ 上的线丛，$s\in H^0(X,L^m)$ 为非零截面，记 $D=(s=0)$。设 $L|_{X\setminus D}$ 是可逆层，且 $(L|_{X\setminus D})^m\cong\cO_{X\setminus D}$。定义 2.49 给出 $m$-重覆盖 $p':Z'\to X\setminus D$，它可延拓为分歧覆盖 $p:Z\to X$。利用 $Z'$ 的代数描述，可以最快地直接得到 $Z$：

% [原文 p.82]
截面 $s$ 可看作层之间的映射 $s:\cO_X\to L^m$。因此 $\bigoplus_{i=0}^{m-1}L^{-i}$ 是 $X$ 上的代数层，其中当 $i+j\ge m$ 时，乘法由
\[
L^{-i}\otimes L^{-j}\cong L^{-i-j}\otimes\cO_X\xrightarrow{\id\otimes s}L^{-i-j}\otimes L^m\cong L^{-i-j+m}
\]
定义。令 $Z=\Spec_X\bigoplus_{i=0}^{m-1}L^{-i}$，并定义投影 $p:Z\to X$。
\end{definition}

$Z$ 的局部描述可立刻得到。取 $x\in X$，选局部坐标 $x_i$，设 $z\in\Gamma(L)$ 为 $x$ 处的局部生成元。$z^m\in\Gamma(L^m)$ 在 $\cO_X$ 中的像为函数 $s(x_1,\dots,x_n)$，它是 $D$ 在 $x$ 处的局部方程。于是在 $x$ 的邻域内，$\bigoplus_{i=0}^{m-1}L^{-i}$ 作为 $\cO_X$-代数由 $z$ 生成，且满足关系 $z^m=s(x_1,\dots,x_n)$。特别地，得到如下简单结果。

\begin{lemma*}{2.51}
$Z$ 为非奇异的充要条件是 $X$ 与 $D$ 都非奇异。
\end{lemma*}

下面需要分歧循环覆盖的如下推广。

\begin{definition}[分歧循环覆盖，一般情形]{2.52}
设 $X$ 为正规簇（或正规解析流形），$L$ 为 $X$ 上秩 1 的无挠层。设 $X^0\subset X$ 为使 $L|_{X^0}$ 局部自由的最大开集。此时 $X\setminus X^0$ 的余维数至少为 2。层 $L^{\otimes i}$ 可以有支集落在 $X\setminus X^0$ 上的挠。

为处理这一情况，引入下述概念。以 $L^{[i]}$ 表示 $L^{\otimes i}$ 的双对偶。$L^{[i]}$ 是 $X$ 上唯一的自反层，满足 $L^{[i]}|_{X^0}\cong(L|_{X^0})^{\otimes i}$。事实上，由此可知：对层 $(L|_{X^0})^{\otimes i}$ 施加的任何操作都自动延拓到这个层。关于循环覆盖，得到如下构造。

设 $s\in H^0(X,L^{[m]})$ 为非零截面，$D=(s=0)$ 为其零点集。$s$ 可看作层之间的同态 $s:\cO_X\to L^{[m]}$。因此 $\bigoplus_{i=0}^{m-1}L^{[-i]}$ 是 $X$ 上的代数层，其中当 $i+j\ge m$ 时，乘法由
\[
L^{[-i]}\otimes L^{[-j]}\cong L^{[-i-j]}\otimes\cO_X\xrightarrow{\id\otimes s}L^{[-i-j]}\otimes L^{[m]}\cong L^{[-i-j+m]}
\]
定义。令 $Z=\Spec_X\bigoplus_{i=0}^{m-1}L^{[-i]}$，并定义投影 $p:Z\to X$。
\end{definition}

\begin{lemmap}{2.53}
记号同上。再假设 $L^{[m]}\cong\cO_X$。此时 $(p^*L)^{[1]}\cong\cO_Z$。
\end{lemmap}

\begin{proof}
$p^*L$ 不必是自反层，故以 $(p^*L)^{[1]}$ 表示其双对偶。如定义 2.49 中所注意到的，$p^*(L|_{X^0})\cong\cO_{p^{-1}(X^0)}$。这一同构自动延拓到整个 $Z$。
\end{proof}

\textbf{步骤 4}\quad 回到原来的设定，假设 $X$ 非奇异且 $L$ 丰富。选取 $m$ 使 $L^m$ 为非常丰富，并取一般截面 $s$。此时 $D=(s=0)$ 是非奇异的非常丰富除子，于是 $Z$ 非奇异，且 $p:Z\to X$ 在 $X\setminus D$ 上是平展的。

% [原文 p.83]
由 Hodge 理论，映射 $\tau:H^i(Z,\C_Z)\to H^i(Z,\cO_Z)$ 是满射。$p$ 的纤维是 0 维的，故高次顺像全为 0，且
\[
p_*\tau:H^i(X,p_*\C_Z)\to H^i(X,p_*\cO_Z)
\]
是满射。

$Z$ 上 $\Z/m\Z$ 的作用把最后这个映射如下分解为特征层的直和。

令 $\xi=e^{2\pi i/m}$。由于 $\Z/m\Z$ 在 $Z$ 上的作用连续，可作如下分解：
\[
p_*\C_Z\cong\bigoplus_{r=0}^{m-1}\C[\xi^r],\qquad p_*\cO_Z\cong\bigoplus_{r=0}^{m-1}L^{-r}.
\]
其中 $\C[\xi^r]$ 是绕除子 $D$ 一周时单值为 $\xi^r$ 的局部系统，对每个 $r$ 都有自然包含 $\C[\xi^r]\to L^{-r}$。$p_*\tau$ 确定对应的特征空间之间的满射。因此，对所有的 $i,r$，映射 $H^i(X,\C[\xi^r])\to H^i(X,L^{-r})$ 都是满射。

\textbf{步骤 5}\quad 对 $0<r<m$，层 $\C[\xi^r]$ 在 $D$ 周围具有非平凡单值，故具有如下简单性质。

\begin{corollary}{2.54}
设 $U\subset X$ 为满足 $U\cap D\ne\varnothing$ 的连通开集。此时 $H^0(U,\C[\xi^r]|_U)=0$。
\end{corollary}

这一性质经由下述明显的结果而被使用。

\begin{lemma}{2.55}
设 $F$ 为拓扑空间 $X$ 上 Abel 群的层，$F_1,F_2\subset F$ 为子层。设 $D\subset X$ 为闭子集。假设下述 2 个条件：
\begin{enumerate}[label=（\arabic*）,leftmargin=2.8em]
\item $F_2|_{X\setminus D}=F|_{X\setminus D}$；
\item 若 $U$ 为连通开集且满足 $U\cap D\ne\varnothing$，则 $H^0(U,F_1|_U)=0$ 成立。
\end{enumerate}
此时 $F_1$ 是 $F_2$ 的子层。
\end{lemma}

% [Tutor Note: 条件 (1) 按页图照录为 $F_2$ 与 $F$ 在 $X\setminus D$ 上一致，缺下标；按上下文与英文底本应为 $F_1$ 与 $F_2$ 一致，读者按补全后的式子理解即可。]

作为其推论，得到：

% [原文 p.84]
\begin{corollaryp}{2.56}
设 $0<r<m$，$M\subset L^{-r}$ 为满足 $M|_{X\setminus D}=L^{-r}|_{X\setminus D}$ 的子层。此时单射 $M\subset L^{-r}$ 使 $\C[\xi^r]\to L^{-r}$ 分解为 $\C[\xi^r]\to M\to L^{-r}$。因此对所有的 $i$，$H^i(X,M)\to H^i(X,L^{-r})$ 都是满射（其中 $0<r<m$）。
\end{corollaryp}

\begin{proof}
前半部分由推论 2.54 与引理 2.55 显然。由此存在映射
\[
H^i(X,\C[\xi^r])\to H^i(X,M)\to H^i(X,L^{-r}).
\]
如上所见，合成映射是满射，因此右侧的映射也是满射。
\end{proof}

\textbf{步骤 6}\quad 准备工作就绪，可以完成论证。令 $M=L^{-r}(-kD)\cong L^{-(r+mk)}$。由推论 2.56，映射 $H^i(X,L^{-(r+mk)})\to H^i(X,L^{-r})$ 对所有 $i,k\ge 0$ 都是满射。取 $r=1$、$k\gg 1$。由对偶性与 Serre 消没定理，左边当 $i<n$ 时等于 0。因此当 $i<n$ 时 $H^i(X,L^{-1})=0$。定理 2.47 的证明至此完成。
\end{proof}

作为小平消没定理的应用，证明除子 $D$ 的 Euler 标数 $\chi(X,\cO_X(D))$ 只依赖于 $D$ 的数值等价类。这是 Riemann–Roch 定理（参看 [Har77, App. A]）的推论；此处对 $\C$ 上的非奇异射影簇给出直接证明。这是后面唯一需要的情形。

\begin{propositionp}{2.57}
设 $X$ 为 $\C$ 上的非奇异射影簇，$D,D'$ 为 $X$ 上的 Cartier 除子，且 $D\equiv D'$。此时
\[
\chi(X,\cO_X(D))=\chi(X,\cO_X(D')).
\]
\end{propositionp}

\begin{proof}
设 $H$ 为 $X$ 上的丰富除子，记 $P(v):=\chi(X,\cO_X(D+vH))$，$P'(u,v):=\chi(X,\cO_X(D+u(D'-D)+vH))$。由定理 1.36，$P$ 与 $P'$ 是 $u,v$ 的多项式。固定任一使 $D+vH-K_X$ 丰富的 $v\gg 1$。此时由 $D'-D\equiv 0$ 及定理 1.18，$D+u(D'-D)+vH-K_X$ 仍为丰富。由定理 2.47，对所有 $u$ 与 $i>0$ 有
\[
H^i(X,\cO_X(D+vH))=0=H^i(X,\cO_X(D+u(D'-D)+vH)).
\]
因此对所有 $u$ 有 $P'(u,v)=h^0(X,\cO_X(D+u(D'-D)+vH))$。由此证明：对每个固定的 $v\gg 1$，$P'(u,v)$ 作为 $u$ 的函数有界。
% [原文 p.85]
这可对 $n=\dim X$ 用归纳法证明如下。取 $m\gg 1$ 使 $(D+vH-mH)\cdot H^{n-1}<0$，并取非奇异除子 $Y\in|mD|$。此时 $H^0(X,\cO_X(D+u(D'-D)+vH-mH))=0$，从而
\[
H^0(X,\cO_X(D+u(D'-D)+vH))\subseteq H^0(Y,\cO_Y((D+u(D'-D)+vH)|_Y)).
\]
有界多项式必为常数，故对 $v\gg 1$，从而对所有 $v$，有 $P'(u,v)=P(v)$。因此
\[
\chi(X,\cO_X(D))=P(0)=P'(1,0)=\chi(X,\cO_X(D')).
\]
\end{proof}

\section{小平消没定理的推广}

\takeaway{本节把消没定理中的丰富性放宽为接近丰富：用截面增长定义大除子，证明大且 nef 的除子可写成丰富 $\Q$-除子与小的有效除子之和，从而把消没定理推广到这类除子。}

本节说明：只要 $L$ 接近丰富，消没定理 2.47 仍然成立。首先明确“接近丰富”这一说法的含义。

\begin{lemmap}{2.58}
设 $X$ 为域上 $n$ 维射影概形，$B$ 为 $X$ 上的 Cartier 除子。此时存在某个数 $C>0$，使得对所有 $k>0$ 有 $h^0(X,\cO_X(kB))\le C\cdot k^n$。
\end{lemmap}

\begin{proof}
取 $X$ 上充分丰富的除子 $H$，使 $H-B$ 线性等价于有效除子。此时 $h^0(X,\cO_X(kB))\le h^0(X,\cO(kH))$，而后者的增长由 Hilbert 多项式给出。若 $X$ 为固有概形，则本引理仍然成立；其证明可经 Chow 引理归到射影情形。
\end{proof}

\begin{definition}{2.59}
设 $X$ 为 $n$ 维固有簇。Cartier 除子 $D$ 若存在 $c>0$，使得对所有 $k\gg 1$ 有 $h^0(X,\cO(kD))>c\cdot k^n$，则称 $D$ 为大的（big）。
\end{definition}

% [Tutor Note: 大性只要求截面增长达到 $k^n$ 的量级，比丰富性弱得多；它对双有理等价不变，因此在极小模型纲领中可以代替丰富性使用。]

大性本质上是丰富性的双有理版本。换言之，若态射 $f:Y\to X$ 双有理，$D$ 为 $X$ 上的 Cartier 除子，则 $D$ 大与 $f^*D$ 大等价。

\begin{lemmap}{2.60}
设 $X$ 为 $n$ 维射影簇，$D$ 为 Cartier 除子。此时以下条件等价：
% [原文 p.86]
\begin{enumerate}[label=（\arabic*）,leftmargin=2.8em]
\item $D$ 大。
\item 存在丰富除子 $A$、有效除子 $E$ 与自然数 $m>0$，使 $mD\sim A+E$ 成立。
\item 对适当的自然数 $m>0$，与线性系 $|mD|$ 相伴的有理映射 $\varphi_{|mD|}$ 是双有理的。
\item 对适当的自然数 $m>0$，$\varphi_{|mD|}$ 的像是 $n$ 维的。
\end{enumerate}
\end{lemmap}

\begin{proof}
注意 (2)$\Rightarrow$(3)$\Rightarrow$(4) 是显然的。为证 (4)$\Rightarrow$(1)，设 $Y:=\varphi_{|D|}(X)\subset\PP^N$ 是 $n$ 维的。由 [Har77, I.7.5]，$Y$ 的 Hilbert 多项式为
\[
h^0(Y,\cO_Y(k))=\frac{\deg Y}{n!}k^n+\text{（低阶项）}.
\]
$\varphi^*_{|D|}$ 引起单射 $H^0(Y,\cO_Y(k))\subset H^0(X,\cO(kD))$，这正说明 (4)$\Rightarrow$(1)。

为证 (1)$\Rightarrow$(2)，设 $D$ 大，$A$ 为丰富且有效的除子。此时有正合列
\[
0\to H^0(X,\cO(kD-A))\to H^0(X,\cO(kD))\to H^0(A,\cO(kD|_A)).
\]
$\dim A=n-1$，故由引理 2.58，$h^0(A,\cO(kD|_A))$ 至多按 $k^{n-1}$ 的量级增长（相差常数倍）。因此对 $k\gg 1$ 有 $H^0(X,\cO(kD-A))\ne 0$。取 $E$ 为 $|kD-A|$ 中的任一有效除子。至此 (1)$\Rightarrow$(2) 得证。
\end{proof}

\begin{propositionp}{2.61}
设 $X$ 为 $n$ 维射影簇，$D$ 为 Cartier 除子。此时以下条件等价：
\begin{enumerate}[label=（\arabic*）,leftmargin=2.8em]
\item $D$ 是 nef 且大的。
\item $D$ 是 nef 的，且满足 $(D^n)>0$。
\item 适当选取有效除子 $E$ 后，$A_k=D-(1/k)E$（$\forall k\gg 1$）为丰富 $\Q$-除子。
\end{enumerate}
若 $X$ 定义在特征 0 的域上，则上述条件还与下述条件等价：
\begin{enumerate}[label=（\arabic*）,leftmargin=2.8em,start=4]
\item 对任一除子 $\Delta\subset X$，存在 $(X,\Delta)$ 的对数奇点消解 $f:Y\to X$ 与 $Y$ 上的有效 snc 除子 $E'$，使得 $A'_k=f^*D-(1/k)E'$（$\forall k\gg 1$）为丰富 $\Q$-除子。
\end{enumerate}
\end{propositionp}
